\section{指定组合的持仓结构与解释}
\subsection{资本权重的可视化}
图\ref{fig:weight_path}展示指定 Global RRP 的全部周度目标权重。上方单独呈现货币市场 ETF，纵轴范围明确标注；下方保留其他 29 只基金的完整时间序列。所有资产使用同一权重分档，不进行逐行归一化，因此相同颜色仍具有相同的资本含义。

\figwide{primary_weights.pdf}{Global RRP 周度目标权重}{fig:weight_path}

热图采用明确的非等宽分档，让小于几个百分点的实际持仓仍可辨识。浅灰表示接近零的权重，蓝色对应较低权重区间，红色对应较高区间。图例直接写明各档上下界，不能将颜色差异解释为等距离的数值变化。分档仅用于显示，不进入优化器，也不改变原始权重。

\subsection{现金与固定收益配置}
主模型平均日度现金类 ETF 权重为 \primaryCashMean，最高为 \primaryCashMax。日利ETF与三只国债、信用债工具的平均合计权重为 \primaryDefensiveMean。该统计包括非调仓日的漂移权重，因此与仅在周度决策日计算的目标权重可以略有差别。

现金类 ETF 的较低波动使其在风险预算参考中具有作用，国债和信用债进一步提高防御性资产占比。该结构降低了组合波动，也限制了增长资产对收益的贡献。

国债和信用债同样不能统一理解为无风险现金。它们承担的利率与信用风险不同，可转债还具有权益相关性。本文将其作为独立工具参与估计，未用单一固定收益收益序列替代。热图可以显示资本在这些工具之间的变化，但不足以单独识别久期或信用因子贡献。

\subsection{资产是否实际使用}
全部 \primaryAssetsUsed{} 只 ETF 都曾在合格期间形成实质持仓。审计采用 $5\times10^{-5}$ 的权重阈值区分数值残差与实际正权重，这一阈值没有作为优化下限。所有资产中，最小的历史峰值权重约为 1.25\%，因而资产参与并非由机器精度附近的微小正数构成。

允许某一期权重为零是优化决策的一部分。只要资产满足准入条件，它仍然参与下一期求解；当前零权重不构成永久剔除。与之相对，上市前或历史不足时的零权重由数据条件决定。热图中两类零权重可能颜色相同，原因需要结合准入记录识别。

\subsection{持仓数量与风险贡献}
资本分散和风险分散仍需分别评价。若多个权益基金高度相关，即使每个权重都不大，也可能共同承担相似市场风险。若某只低波动资产权重大，其资本占比与风险贡献也可能相差较大。全部资产曾参与只能说明机会集合获得使用，不能说明每一期风险均匀分配。

参考组合的贡献误差与最终组合的贡献误差应分别读取。前者验证风险预算参考的构造，后者反映收益风险松弛造成的偏离。最终问题并未要求贡献完全相等，所以更大的最终误差不能自动作为求解失败；应结合预算、目标值和预测风险检查。

\subsection{持仓变化与收益解释}
对某一期资产收益进行解释，可以从 $w_{i,t}r_{i,t}$ 的日度贡献出发。其跨日简单求和属于算术贡献，不能不加说明地称为累计复利收益的精确分解。多期归因还需要指定复利链接方式和费用分摊。本文不以热图颜色代替严格的多期收益归因。

同理，某资产在最大回撤阶段权重较高，也不能单凭相关时间关系断言它导致全部回撤。权重会随周度决策改变，其他资产收益和交易成本也会同时进入组合。本文采用净值、回撤和持仓的相互参照来形成解释，但不制造未经核算的因果百分比。

\subsection{仓库记录与论文展示}
完整周度记录由三类 CSV 组成。权重矩阵支持绘图和逐期比较，持仓明细支持核对交易前权重、目标及买卖金额，周收益摘要区分自然周与持有期收益。三个表共享决策日期和资产映射，可以与日收益表相互校验。
